\section{Balanced phase and quantum controllability}
\label{sec:balanced-control}
This section proves an exact finite-dimensional Lie-algebra statement.
Controllability allows varying controls; it is not a local-statistics theorem
for one fixed kicked propagator. The underlying Lie-algebra mechanism is a
standard type of quantum-control criterion, and priority is not asserted.

Let $N\ge3$ be odd and use modular-centred Weyl operators, so
\[
 W(v)W(w)=e_N(2^{-1}\det(v,w))W(v+w),\qquad W(v)^N=I,
\]
where $2^{-1}$ is the inverse modulo $N$. Suppose $\det(v,w)$ is a unit
modulo $N$, and define
\[
 H_0=(1+i)W(v)+(1-i)W(-v),\qquad
 H_1=cW(w)+\overline cW(-w),\quad c\ne0.
\]

\begin{theorem}[An explicit generating pair]\label{thm:balanced-su}
The real Lie algebra generated by $iH_0,iH_1$ is
$\mathfrak{su}(N)$. Composite odd dimensions are included.
\end{theorem}
\begin{proof}
Put $\zeta=e^{2\pi i/N}$ and $d=\det(v,w)$. Weyl commutation gives a basis
in which $W(v)$ has simple eigenvalues $\zeta^j$ and $W(w)$ maps the $j$th
basis vector to the $(j+d)$th with a nonzero weight. Thus $H_1$ has a
connected cyclic edge graph. The eigenvalues of $H_0$ are
\[
 E_j=2\sqrt2\cos(2\pi j/N+\pi/4).
\]
They are distinct: a nontrivial cosine equality would require
$j+k+N/4$ to be an integer multiple of $N$. The edge gaps are
\[
 \Delta_j=-4\sqrt2\sin(\pi d/N)
               \sin(2\pi j/N+\pi d/N+\pi/4).
\]
Their squares are pairwise distinct. Equality of squared sines implies
either $2(j-k)=0\pmod N$, hence $j=k$, or
$4(j+k+d)=(2a-1)N$, impossible because the right side is odd.
The gaps are nonzero because the $E_j$ are distinct.

On the skew-Hermitian subspace of edge $j$, $\operatorname{ad}(iH_0)^2$
has eigenvalue $-\Delta_j^2$. Real polynomial interpolation isolates each
edge component of $iH_1$. Commutation with $iH_0$ gives its other real
quadrature. Their commutators give diagonal differences, and commutators
along paths in the connected cycle give all off-diagonal generators.
These span $\mathfrak{su}(N)$.
\end{proof}

Let $U_NW(v)U_N^*=W(Av)$ and $D=\det(v,Av)\ne0$. For every odd $N$
coprime to $D$, the theorem applies to $H_0,U_NH_0U_N^*$.
For the Pell matrix $A=\left(\begin{smallmatrix}3&2\\4&3\end{smallmatrix}\right)$
and $v=(1,0)$, $D=4$, so every odd dimension is admissible. This matrix is
identity modulo two; the physical integer lift has the same covariance
along this orbit.

\begin{theorem}[One parameter in all dimensions]\label{thm:simultaneous-control}
Let $K_N$ be arbitrary specified traceless Hermitian backgrounds. For each
admissible $N$, a nonzero real polynomial of degree at most $(N^2-1)^2$
certifies that
\[
 \operatorname{Lie}_{\mathbb R}
 \bigl(i(K_N+\varepsilon H_0),
       iU_N(K_N+\varepsilon H_0)U_N^*\bigr)=\mathfrak{su}(N)
\]
outside its roots. Hence one real $\varepsilon$, outside a countable
exceptional set, works simultaneously in every admissible dimension.
Such choices exist arbitrarily close to zero.
\end{theorem}
\begin{proof}
Write $D_N=N^2-1$. The spaces spanned by left-nested Lie words of length at
most $L$ strictly increase until the generated algebra is reached. If one
step fails to increase, the space is invariant under both generator
adjoints and no later step can increase it. Thus at the pure $H_0$ pair
choose $D_N$ independent homogeneous words of lengths at most $D_N$.
Evaluate these words on the perturbed pair. Their coordinate determinant in
a fixed real basis of $\mathfrak{su}(N)$ is polynomial in $\varepsilon$.
Its leading coefficient is the nonzero determinant of the pure-pair words;
its degree is at most the sum of their lengths, hence at most $D_N^2$.
Take the countable union of the finite root sets. Roots of this particular
certificate need not be genuinely uncontrollable parameters.
\end{proof}

For a real even classical background, the balanced harmonic
\[
 b=(1+i)e_v+(1-i)e_{-v}
\]
also gives the full classical Riesz property of
Theorem~\ref{thm:full-riesz}. Its even part adds
$\varepsilon(e_v+e_{-v})$ and its sine part adds
$\varepsilon i(e_v-e_{-v})$. Hence
\[
 c_\infty(g+\varepsilon b)
 \ge\frac{4\pi^2\varepsilon^4D^2}{(\operatorname{tr}A)^2-4}.
\]
For the Pell example this is $2\pi^2\varepsilon^4$.
Thus an arbitrarily small fixed harmonic perturbation can simultaneously
regularize the entire classical resonance lattice and give full quantum
control in every admissible odd dimension, outside the stated countable set.

There is no dimension-uniform synthesis-time estimate here. Nor is there a
mixing-rate estimate for random controls, or a Heisenberg-scale estimate for
a fixed propagator. In fact if an unperturbed $U_N$ has scalar $K_N$th power,
the lag commutators
\[
 \frac{iN}{2}[G,U_N^rGU_N^{-r}]
\]
are $K_N$-periodic and vanish at $r=K_N$. Any Gram matrix of more than
$K_N-1$ such lag observables is singular, regardless of nonzero column
normalizations. A classical Gram block bounded below by $cI$ cannot be
approximated in operator norm within $c$ by that singular quantum block.
This is an obstruction to one transfer mechanism, not a conclusion about
an actually perturbed propagator whose exact period has been destroyed.
